\section{Temporary membership questions and their regret cost}
\label{sec:tradeoff}

\subsection{A repeated-update family}

We now keep the reward environment, arrival schedule, endpoint, and reporting
parameters fixed when comparing two area budgets. Choose integers $J\ge1$,
$M\ge\max\{2,J\}$ and
\begin{equation}
 0<\varepsilon\le1/8,\quad 0<\delta\le1/16,\quad
 Jv\le\varepsilon/4,\quad 0<d\le v/32,\quad\eta=d/2.
 \label{eq:family_parameters}
\end{equation}
At $a_e=1+eH$, $e=0,\ldots,J$, a reference with mean $b_e=1/2+ev$
arrives. For each $e<J$, the same batch contains $M$ candidates with Bernoulli
means
\begin{equation}
 \mu_{e,k}=b_e-\varepsilon+\sigma_{e,k}d,
           \qquad \sigma_{e,k}\in\{-1,+1\}.
 \label{eq:family_means}
\end{equation}
The last batch contains only its reference, and $T=(J+1)H$. The family
$\Theta$ includes every sign configuration and every within-batch permutation.
The policy is not told which arm is the reference. Lower bounds may disclose
those identities and the schedule. All arms remain available.

In its birth window a candidate is good exactly when its sign is positive.
Its regret gap is approximately $\varepsilon$, but its distance to the
membership boundary is $d$. At the next update both signs become bad, with
boundary distance at least $v-d$. Every reference remains good. Thus the final
good set consists only of references, regardless of all the candidate signs.
The statistical tasks that distinguish the environments are temporarily
relevant and can expire while their arms remain available.

\subsection{Area forces completion of enough temporary tasks}

The lower-bound ingredient permits arbitrary interleaving and stopping. For
$m$ independent uniform signs with observation laws $P_i^+,P_i^-$, put
\[
 \kappa=\max_i\max\{D(P_i^+\Vert P_i^-),D(P_i^-\Vert P_i^+)\}
                                                    \in(0,\infty).
\]
If a bounded stopping time has at least half the signs determined with prior
probability at least $1/2$, and all configurations have familywise report
error at most $\delta\le1/16$, then
\begin{equation}
 \overline\E\sum_{i=1}^mN_i\ge
                 \frac{m}{1024\kappa}\log\frac m\delta.
 \label{eq:partial_recovery}
\end{equation}
Other observations must not encode these signs. The proof factors the posterior
after the sampling and stopping kernels cancel, and combines a simultaneous
recovery bound with coordinatewise change of measure
\citep{kaufmann2016}. It does not assert independence of the random sample
counts. The logarithmic simultaneous-recovery scale is established statistical
structure; its role here is to price partially completed, expiring tasks.

For clarity, a heterogeneous version gives an explicit lower allocation
constraint. Let a compatible batch $e$ have $M_e$ candidates, window length
$H_e$, boundary distance $d_e>\eta$, and two-way KL bounded by $\kappa_e$.
Put $w_e=M_eH_e$, $V=\sum_ew_e$. Define
\[
 \Phi_\delta(u)=
 \begin{cases}\kl(1-\delta-u,\delta),&u<1-2\delta,\\0,&u\ge1-2\delta,
 \end{cases}
\]
\[
 g_{m,\delta}(p)=\max\left\{0,
 \frac{(p-2\delta)_+}{2}\left[\log\frac m{4\log2}\right]_+-1,
 \Phi_\delta(1-p/2)\right\}.
\]
If a policy satisfies the joint area requirement with budget $a$, then
\begin{equation}
 \sup_\sigma\E_\sigma R_T\ge
 \inf_{\substack{0\le p_e\le1\\\sum_ew_ep_e\ge(1-\delta)V-2a}}
 \sum_e\frac{(\varepsilon-d_e)M_e}{\kappa_e}g_{M_e,\delta}(p_e).
 \label{eq:allocation_lower}
\end{equation}
Here $p_e$ is the probability of determining at least half of batch $e$ before
its window ends. A window that never does so contributes at least $w_e/2$
area. This gives the constraint directly from the area success event, with its
failure probability retained. The objective then uses
\eqref{eq:partial_recovery} under one compatible product prior. The policy can
choose where to defer; no separate requirement forces every window to finish.
The relaxation need not be attainable in general.

\subsection{Reusable coarse evidence and a matched rapid endpoint}

CP-UCB-AA can use a predetermined precision $\rho_e$ for each arriving batch.
Let $S_e=\sum_{i:\tau_i=a_e}q_i(\rho_e)$ be the total checkpoint work of
batch $e$, and set $M_{-1}=M_J=0$. In a heterogeneous repeated-update family,
suppose
\begin{equation}
 H_e\ge2S_e,\qquad b_J-b_0\le\varepsilon/4,\qquad
 d_e+\rho_e+\rho_{e+1}<b_{e+1}-b_e,\qquad
 2\max_e\rho_e\le\varepsilon/2.
 \label{eq:reuse_conditions}
\end{equation}
These are conditions on the performance family, not future inputs supplied to
the policy. Write $\mathcal F=\{e<J:2\rho_e<d_e\}$.

\begin{proposition}[Shared update from lifetime evidence]
\label{thm:reuse}
On the checkpoint coverage event, under \eqref{eq:reuse_conditions},
\begin{equation}
 \mathcal A_T(\eta)\le
 \sum_{e\notin\mathcal F}M_eH_e+
 2\sum_{e\in\mathcal F}M_eS_e+
 2\sum_{e=0}^J(e+1+M_{e-1})S_e.
 \label{eq:reuse_area}
\end{equation}
When the next batch completes its checkpoints, every preceding-batch candidate
can be excluded using its saved interval and the new reference lower bound,
without another observation of that old candidate for certification.
\end{proposition}

The last term charges only temporarily affected references and the preceding
candidate cohort, rather than all historical candidates at each update.
The conclusion follows because the new lower bound exceeds every preceding
candidate upper bound by more than $\varepsilon$. The condition
$H_e\ge2S_e$ ensures that all its checkpoint observations occur in its birth
window. The corresponding unconditional reward bound, including excess UCB
pulls, is given in the supplement.

For the homogeneous family \eqref{eq:family_parameters}--\eqref{eq:family_means},
set
\begin{gather}
 K=J(M+1)+1,\quad L_K=\log\frac{4K(K+1)}\delta,\quad
 \rho_f=d/4,\quad\rho_c=v/8,\notag\\
 q_f=\lceil32L_K/d^2\rceil,\quad q_c=\lceil128L_K/v^2\rceil,\quad
 H=128(M+J+1)q_f,\quad V=JMH.
 \label{eq:chosen_horizon}
\end{gather}
The quantities $q_f,q_c$ are uniform analysis bounds on the actual armwise
quotas. Neither final arm count nor future times are policy inputs. Fine
precision can be set from the task input as $\rho_f=\eta/2$; the coarse
precision is a separate preselected configuration. Automatic adaptation to
an unknown improvement scale is not asserted.

\begin{theorem}[Rapid and deferred maintenance at the same endpoint]
\label{thm:tradeoff}
On the preceding family and horizon,
\begin{equation}
 2^{-20}\frac{\varepsilon JM}{d^2}\log\frac M\delta
 \le\mathfrak R_\Theta(V/8)\le
 2^{18}\frac{\varepsilon JM}{d^2}\log\frac M\delta.
 \label{eq:rapid_matched}
\end{equation}
The coarse-precision CP policy has joint label-and-area success probability
at least $1-\delta$ with area budget $9V/8$, and
\begin{equation}
 \sup_{\theta\in\Theta}\E_\theta R_T\le
 2\varepsilon JM q_c+
 64\left(\frac{JM}{\varepsilon}+\frac{J^2}{v}\right)\log(2T)+1.
 \label{eq:coarse_regret}
\end{equation}
Consequently, at fixed legal $J,M,\varepsilon,v,\delta$ and $d\downarrow0$,
with $\eta=d/2$ and the horizon in \eqref{eq:chosen_horizon},
$\mathfrak R_\Theta(V/8)=\Theta(d^{-2})$, whereas
$\mathfrak R_\Theta(9V/8)=O(\log(1/d))$.
\end{theorem}

For the lower bound, area at most $V/8$ forces half-completion in at least
$3J/4$ windows on success. Enough windows then require
\eqref{eq:partial_recovery}, and each candidate observation costs at least
$\varepsilon-d$. Distinct batches' samples are charged once. For the upper
bound, \eqref{eq:reuse_area} gives the appropriate area budgets. At the specified
horizon the deterministic bound $R_T\le2\varepsilon T$ suffices for the rapid
order; it would not justify the same conclusion at arbitrarily longer windows.
The coarse bound instead counts checkpoint observations and excess ordinary
UCB pulls separately.

The two budgets differ by a constant factor, but they allow different service
levels. The theorem does not claim that the coarse policy meets the rapid
budget, that its regret upper bound is matched, or that all intermediate budgets
are characterized. A preselected mixture of fine and coarse batches gives
additional attainable points, as recorded in the supplement.

\subsection{A probability lower bound under the same area requirement}

An expectation lower bound alone could, in principle, reflect unusual sample
paths. The following result gives a stronger conclusion without replacing the
cumulative area requirement by deadlines.

\begin{theorem}[A nonvanishing probability of the rapid-maintenance cost]
\label{thm:quantile_main}
In the repeated-update family, suppose a reliable policy satisfies
\[
 \Prob_\sigma\{E_{\rm lab},\mathcal A_T(\eta)\le V/8\}\ge1-\delta
\]
for every sign configuration. Then one fixed configuration satisfies
\begin{equation}
 \Prob_\sigma\left\{R_T\ge
 2^{-20}\frac{\varepsilon JM}{d^2}\log\frac M\delta\right\}
 \ge\frac{11}{80}.
 \label{eq:quantile_main}
\end{equation}
\end{theorem}

The proof stops each window at half-completion, its endpoint, or the first
candidate sample beyond a fixed budget. Applying the partial-recovery lemma
to this stopped experiment bounds the probability of a cheap completion by
$1/2$. Low total area requires many completions; a count of cheap completions
then leaves enough expensive windows with probability at least
$1/5-\delta\ge11/80$. Their actual observations give
\eqref{eq:quantile_main}. No independence between window-completion events is
needed. Since $11/80>\delta$, the rapid endpoint in
\eqref{eq:rapid_matched} also has matching order when regret and area are
bounded on one event of probability at least $1-\delta$. The full stopping
argument and a compatible multiple-window area truncation are proved in the
supplement.
